\chapter{The Interdependent Task Allocation Problem}\label{chap:task_allocation}

\subsection{Centralized ITAP}

The scheduling models of the previous chapter provide a useful language for allocation under uncertainty, partial revelation, and precedence structure. However, the research-allocation problems motivating this thesis are not fully captured by classical scheduling alone. In mathematical development, software verification, and related domains, tasks do not merely wait to be processed: they may reveal further tasks, alter the effective difficulty of future work, admit multiple alternative solution routes, and interact with heterogeneous agents in strongly state-dependent ways. The purpose of this section is to formalize a centralized benchmark model that captures these phenomena while retaining the operational flavor of stochastic online scheduling. This model is the \emph{Interdependent Task Allocation Problem} (ITAP) in its centralized regime.

\subsubsection{Motivation}

The motivating picture is the one already suggested in the introduction: major intellectual achievements rarely arise from a single isolated task solved in a vacuum. Rather, they emerge from the coordinated resolution of many interdependent tasks, some visible from the outset and others only becoming apparent once earlier work has been completed. In mathematics, this is especially clear. A target theorem may depend on a large body of intermediate lemmas, conceptual reorganizations, and auxiliary infrastructure; some of these dependencies are known in advance, while others are only discovered as the work progresses. More broadly, the same pattern appears in theorem proving, formal verification, experimental science, and engineering design: progress is not merely the execution of a fixed list of jobs, but the navigation of an evolving dependency structure over tasks whose values are largely downstream.

Classical scheduling captures part of this story. It allows one to reason about heterogeneous resources, release dates, precedence constraints, preemption, and stochastic processing times. Yet it still treats jobs as largely exogenous objects whose significance is intrinsic and whose role in the future of the system is secondary. By contrast, in the settings motivating this thesis, the importance of a task is often determined less by its immediate completion and more by what its completion reveals, enables, or disambiguates. Solving a task may unlock new tasks, certify that some alternative routes are unnecessary, change the effective tractability of other tasks, or expose previously latent structure in the dependency graph itself.

This motivates viewing research and related forms of problem-solving as an \emph{interdependent} allocation problem. We imagine a central planner allocating a collection of heterogeneous agents across a pool of partially revealed tasks. Each task may admit several possible modes of attack, several agents may collaborate on the same task, and the time required to complete a task may depend on both the current state of completed work and latent environmental factors that are not fully visible to the planner. At the same time, completing tasks changes the future by revealing new tasks and by modifying the set of feasible or promising continuations.

The centralized regime is the analytically clean starting point for this study. In this regime there is a single planner, a single objective, and no private information or strategic behavior among agents. Thus the centralized ITAP should be understood as a controlled, belief-augmented generalization of the stochastic online scheduling models introduced earlier. It serves as both a benchmark model in its own right and as the conceptual bridge to the decentralized regime considered later, where private information and strategic interaction can no longer be ignored.

\subsubsection{Formal Definition}

We now define the centralized ITAP as a particular belief-augmented generalized stochastic online scheduling instance. Throughout this section, we work in discrete time \(t\in\mathbb N_0\), in keeping with the discrete-time focus adopted earlier.

\begin{definition}[Centralized Interdependent Task Allocation Problem]
A \emph{centralized Interdependent Task Allocation Problem} (centralized ITAP) instance is a tuple
\[
\mathfrak I=
\bigl(
\mathcal A,\,
\mathcal T^\ast,\,
\{\mathcal O_\tau\}_{\tau\in\mathcal T^\ast},\,
\{\mathcal W_\tau\}_{\tau\in\mathcal T^\ast},\,
(\Omega,\mathcal F,\mathbb P),\,
(\Theta,\mathcal T,\theta),\,
\mathfrak R,\,
\{L_{\tau,o}\}_{\tau\in\mathcal T^\ast,\ o\in\mathcal O_\tau},\,
\{\rho_{\tau,o}\}_{\tau\in\mathcal T^\ast,\ o\in\mathcal O_\tau},\,
\mathfrak Z,\,
\Gamma
\bigr),
\]
consisting of the following data.

\begin{enumerate}
    \item \textbf{Agents.}  
    A finite set of agents
    \[
    \mathcal A=\{a_1,\dots,a_m\}.
    \]
    These are the centralized analogue of the machines in the preceding scheduling models, but we use the term ``agent'' to emphasize that they may represent mathematicians, theorem provers, verification engines, experimental resources, or other heterogeneous workers.

    \item \textbf{Latent task set.}  
    A finite or countable latent task set
    \[
    \mathcal T^\ast.
    \]
    Elements \(\tau\in\mathcal T^\ast\) are the tasks that may in principle arise during the execution of the system. The full latent task set need not be visible at time \(0\).

    \item \textbf{Modes of attack / operations.}  
    For each task \(\tau\in\mathcal T^\ast\), a finite nonempty set
    \[
    \mathcal O_\tau
    \]
    of operations or modes of attack for task \(\tau\). Intuitively, these encode different ways of attempting the task. In the simplest case one may take \(|\mathcal O_\tau|=1\) for every \(\tau\), but the general formalism permits multiple alternative routes.

    \item \textbf{Completion semantics.}  
    For each task \(\tau\in\mathcal T^\ast\), a nonempty family
    \[
    \mathcal W_\tau\subseteq 2^{\mathcal O_\tau}\setminus\{\varnothing\}.
    \]
    A task \(\tau\) is considered complete once the set of completed operations for \(\tau\) contains some \(W\in\mathcal W_\tau\). Thus:
    \begin{itemize}
        \item the classical ``AND'' case is recovered by taking \(\mathcal W_\tau=\{\mathcal O_\tau\}\);
        \item the classical ``OR'' case is recovered by taking
        \[
        \mathcal W_\tau=\bigl\{\{o\}:o\in\mathcal O_\tau\bigr\}.
        \]
    \end{itemize}

    \item \textbf{Latent environment and prior.}  
    A probability space
    \[
    (\Omega,\mathcal F,\mathbb P)
    \]
    together with a measurable latent parameter space \((\Theta,\mathcal T)\) and a measurable random element
    \[
    \theta:\Omega\to\Theta.
    \]
    The latent parameter \(\theta(\omega)\) encodes all hidden environmental features relevant to the instance, such as latent difficulty classes, latent dependency information, latent tractability of operations, or any other hidden primitive data. The prior law of \(\theta\) under \(\mathbb P\) is the planner's ex ante belief.

    \item \textbf{Task revelation operator.}  
    A measurable map
    \[
    \mathfrak R:2^{\mathcal T^\ast}\times \Theta \to 2^{\mathcal T^\ast}
    \]
    such that for every \(\vartheta\in\Theta\):
    \begin{itemize}
        \item \(\mathfrak R(C,\vartheta)\supseteq C\) for all \(C\subseteq\mathcal T^\ast\);
        \item if \(C\subseteq C'\), then
        \[
        \mathfrak R(C,\vartheta)\subseteq \mathfrak R(C',\vartheta).
        \]
    \end{itemize}
    Thus \(\mathfrak R(C,\vartheta)\) is the set of tasks revealed once the currently completed task set is \(C\) and the latent environment is \(\vartheta\). The monotonicity assumption formalizes the idea that completing more tasks cannot erase previously revealed tasks.

    We write
    \[
    \mathcal T_0^\vartheta := \mathfrak R(\varnothing,\vartheta)
    \]
    for the initially revealed task set under latent environment \(\vartheta\).

    \item \textbf{Latent workloads.}  
    For each task \(\tau\in\mathcal T^\ast\) and operation \(o\in\mathcal O_\tau\), a latent workload random variable
    \[
    L_{\tau,o}:\Omega\to (0,\infty].
    \]
    Here \(L_{\tau,o}(\omega)\) is the total amount of effective work required to complete operation \(o\) of task \(\tau\) in scenario \(\omega\). The value \(+\infty\) is allowed, thereby permitting intractable operations.

    \item \textbf{State-dependent productivity / progress rule.}  
    For each task \(\tau\in\mathcal T^\ast\) and operation \(o\in\mathcal O_\tau\), a measurable map
    \[
    \rho_{\tau,o}:2^{\mathcal A}\times 2^{\mathcal T^\ast}\times \Theta \to [0,\infty),
    \]
    where
    \[
    \rho_{\tau,o}(S,C,\vartheta)
    \]
    is the instantaneous effective work rate on operation \(o\) when the set of agents \(S\subseteq \mathcal A\) is simultaneously assigned to \(o\), the completed task set is \(C\), and the latent environment is \(\vartheta\).

    This map encodes:
    \begin{itemize}
        \item heterogeneity of agent capabilities;
        \item redundancy or collaboration, since \(|S|\) may exceed \(1\);
        \item state dependence of task difficulty through the dependence on \(C\).
    \end{itemize}

    We require \(\rho_{\tau,o}(\varnothing,C,\vartheta)=0\) for all \(\tau,o,C,\vartheta\).

    \item \textbf{Visible signals / observation rule.}  
    A family \(\mathfrak Z\) of measurable signals specifying what the centralized planner observes when tasks are revealed and when operations complete. Concretely, this includes at least:
    \begin{itemize}
        \item for each task \(\tau\), a revelation signal describing the metadata that become visible when \(\tau\) enters the revealed task set;
        \item for each operation \((\tau,o)\), a completion signal describing the information that becomes visible when \(o\) completes.
    \end{itemize}
    In particular, the formalism allows that the planner may observe only partial information about the latent workload law or dependency structure at task revelation time, and must therefore act on a posterior belief rather than on full distributional knowledge.

    \item \textbf{Objective functional.}  
    A measurable performance functional
    \[
    \Gamma
    \]
    acting on complete trajectories of the controlled process. Depending on the application, \(\Gamma\) may minimize time to complete a target set, maximize expected cumulative value of completed tasks, minimize weighted completion times, or optimize any other measurable functional of the task-allocation trajectory.
\end{enumerate}
\end{definition}

\begin{remark}[Interpretation of the revelation operator]
The revelation operator \(\mathfrak R\) is the formal device by which the centralized ITAP captures recursive task generation. A task need not be visible from the outset. Instead, it may become visible only after certain other tasks have been completed. One may equivalently think of \(\mathfrak R\) as being induced by a latent dependency hypergraph, but for the centralized ITAP it is convenient to work directly with the monotone revelation operator, since it captures both ordinary precedence constraints and more general completion-triggered task revelation.
\end{remark}

\begin{remark}[Interpretation of state-dependent productivity]
The maps \(\rho_{\tau,o}\) encode the fact that the effective difficulty of a task may depend on the currently completed set. This is the key point at which the ITAP departs from ordinary scheduling. Completing one task may simplify another, render a previously ambiguous path straightforward, or make an alternative route unnecessary. Thus task difficulty is not treated as a fixed exogenous label, but as a state-dependent quantity.
\end{remark}

\begin{remark}[Beliefs and partial visibility]
The latent parameter \(\theta\) and the observation family \(\mathfrak Z\) together encode the planner's partial knowledge. At any time \(t\), the centralized planner need not know the full latent environment, nor even the full laws of the unresolved tasks. Instead, it acts on the posterior belief induced by the revealed history. Classical stochastic online scheduling is recovered as the special case in which, upon task revelation, the relevant processing-time distribution is fully visible.
\end{remark}

We now define the induced state, action, and policy objects for a fixed centralized ITAP instance.

\begin{definition}[Centralized ITAP state]
Fix a centralized ITAP instance \(\mathfrak I\). A state at time \(t\) is a tuple
\[
X_t=
\bigl(
\mathcal T_t,\,
\mathcal C_t,\,
\{ \mathcal C_{\tau,t}^{\mathrm{op}}\}_{\tau\in\mathcal T_t},\,
B_t,\,
Y_t,\,
\zeta_t,\,
\Pi_t
\bigr),
\]
where:
\begin{itemize}
    \item \(\mathcal T_t\subseteq\mathcal T^\ast\) is the set of tasks revealed by time \(t\);
    \item \(\mathcal C_t\subseteq\mathcal T_t\) is the set of tasks completed by time \(t\);
    \item for each revealed task \(\tau\in\mathcal T_t\),
    \[
    \mathcal C_{\tau,t}^{\mathrm{op}}\subseteq \mathcal O_\tau
    \]
    is the set of completed operations of \(\tau\) by time \(t\);
    \item
    \[
    B_t:\mathcal A\to \Bigl(\bigsqcup_{\tau\in\mathcal T_t}\mathcal O_\tau\Bigr)\cup\{\varnothing\}
    \]
    records the current centralized assignment of each agent, with \(B_t(a)=\varnothing\) meaning that agent \(a\) is idle;
    \item
    \[
    Y_t(\tau,o)\in[0,\infty)
    \]
    is the cumulative effective work accrued on operation \(o\in\mathcal O_\tau\) up to time \(t\);
    \item \(\zeta_t\) records all scheduler-visible metadata and signals revealed by time \(t\);
    \item
    \[
    \Pi_t := \operatorname{Law}(\theta\mid H_t)
    \]
    is the posterior belief over latent environments given the revealed history \(H_t\) up to time \(t\).
\end{itemize}

These data are required to satisfy:
\begin{enumerate}
    \item \(\mathcal T_t=\mathfrak R(\mathcal C_t,\theta)\) at the latent level, with only the visible portion represented in \(\zeta_t\);
    \item a task \(\tau\in\mathcal T_t\) is complete at time \(t\) if and only if there exists \(W\in\mathcal W_\tau\) such that
    \[
    W\subseteq \mathcal C_{\tau,t}^{\mathrm{op}};
    \]
    \item an operation \(o\in\mathcal O_\tau\) is complete once its cumulative effective work reaches its latent workload:
    \[
    Y_t(\tau,o)\ge L_{\tau,o}(\omega).
    \]
\end{enumerate}
\end{definition}

\begin{definition}[Centralized admissible actions]
Fix a centralized ITAP instance \(\mathfrak I\) and a state \(x\). An admissible action at \(x\) is a map
\[
u:\mathcal A\to \Bigl(\bigsqcup_{\tau\in\mathcal T^x}\mathcal O_\tau\Bigr)\cup\{\varnothing\}
\]
such that:
\begin{itemize}
    \item if \(u(a)=(\tau,o)\), then \(\tau\in\mathcal T^x\setminus\mathcal C^x\);
    \item the operation \(o\in\mathcal O_\tau\) has not already been completed;
    \item any additional revealed feasibility constraints encoded in \(x\) are respected.
\end{itemize}

Multiple agents may be assigned to the same operation simultaneously. Thus centralized actions are global assignment vectors, chosen by a single planner, and not independent local actions of the agents.
\end{definition}

\begin{definition}[Centralized nonanticipative policy]
For each \(t\ge 0\), let \(\mathsf H_t^{\mathfrak I}\) denote the measurable space of feasible histories of centralized ITAP states up to time \(t\). A \emph{deterministic centralized nonanticipative policy} is a family of measurable maps
\[
\varphi_t:\mathsf H_t^{\mathfrak I}\to \mathsf U_{\mathfrak I},\qquad t\in\mathbb N_0,
\]
such that
\[
\varphi_t(h)\in \mathcal U_{\mathfrak I}(x_t(h))
\qquad\text{for every }h\in \mathsf H_t^{\mathfrak I},
\]
where \(x_t(h)\) is the terminal state of the history \(h\), \(\mathsf U_{\mathfrak I}\) is the global action space, and \(\mathcal U_{\mathfrak I}(x)\) is the admissible-action set at state \(x\).

A \emph{randomized centralized nonanticipative policy} is a family of stochastic kernels
\[
K_t(\cdot\mid h)\in\mathcal P(\mathsf U_{\mathfrak I})
\]
supported on \(\mathcal U_{\mathfrak I}(x_t(h))\).

Thus centralized policies are precisely the global, history-dependent allocation rules by which a single planner chooses, at each time, how all agents are assigned across the currently revealed tasks.
\end{definition}

\begin{remark}[Centralized ITAP as a specialization of generalized SOS]
The centralized ITAP is a particular belief-augmented generalized stochastic online scheduling instance in which:
\begin{itemize}
    \item jobs are reinterpreted as tasks;
    \item machines are reinterpreted as agents;
    \item the revealed task set is generated endogenously through the revelation operator \(\mathfrak R\);
    \item state-dependent processing laws are encoded by the productivity maps \(\rho_{\tau,o}\);
    \item multiple agents may redundantly collaborate on a single operation;
    \item and task completion may follow generalized OR/AND semantics through the families \(\mathcal W_\tau\).
\end{itemize}
This is the sense in which the centralized ITAP extends the earlier online stochastic scheduling formalism while remaining within its general conceptual envelope.
\end{remark}

\begin{remark}[Centralization]
The qualifier ``centralized'' is important. All allocations are chosen by a single planner with a single objective functional \(\Gamma\). Agents do not have private information, private utilities, or strategic incentives in this model. Those features will be introduced only in the decentralized regime.
\end{remark}

\subsubsection{Simple Examples}


\subsection{Approximation Theory}

\subsection{Complexity Theory}

\subsection{Implications and Limitations}


